% Figure 12.3: overlapping divergences in the electron self-energy.
% Source: physical PDF p. 535 (printed p. 512).
\begingroup
\renewcommand{\thefigure}{12.\arabic{figure}}
\setcounter{figure}{2}
\begin{figure}[tbp]
\centering
\begin{tikzpicture}[
  electron/.style={line width=0.7pt,
    postaction={decorate},decoration={markings,
      mark=at position 0.52 with {\arrow{Stealth[length=4.6pt,width=3.5pt]}}}},
  photon/.style={line width=0.65pt,decorate,
    decoration={snake,amplitude=1.0pt,segment length=4.8pt}},
  callout/.style={-{Stealth[length=4.6pt,width=3.5pt]},line width=0.65pt},
  every node/.style={font=\small}
]
  % Overlapping two-loop electron self-energy.
  \begin{scope}[xshift=-3.25cm,yshift=1.35cm]
    \draw[line width=0.7pt] (-3.00,0) -- (-1.45,0);
    \draw[line width=0.7pt] (1.45,0) -- (3.00,0);
    \draw[electron] (-1.45,0) .. controls (-1.18,0.68) and (-0.52,0.88) .. (0,0.88);
    \draw[photon] (-1.45,0) .. controls (-1.18,-0.68) and (-0.52,-0.88) .. (0,-0.88);
    \draw[line width=0.7pt] (0,0.88) -- (0,-0.88);
    \draw[photon] (0,0.88) .. controls (0.52,0.88) and (1.18,0.68) .. (1.45,0);
    \draw[electron] (0,-0.88) .. controls (0.52,-0.88) and (1.18,-0.68) .. (1.45,0);
  \end{scope}
  % Vertex counterterm at the left end of the loop.
  \begin{scope}[xshift=3.25cm,yshift=1.35cm]
    \draw[line width=0.7pt] (-3.00,0) -- (-1.45,0);
    \draw[line width=0.7pt] (1.45,0) -- (3.00,0);
    \draw[photon] (-1.45,0) arc[start angle=180,end angle=0,x radius=1.45cm,y radius=0.88cm];
    \draw[electron] (-1.45,0) arc[start angle=180,end angle=360,x radius=1.45cm,y radius=0.88cm];
    \draw[line width=0.8pt] (-1.57,-0.12)--(-1.33,0.12);
    \draw[line width=0.8pt] (-1.57,0.12)--(-1.33,-0.12);
    \draw[callout] (-2.05,-1.15) -- (-1.62,-0.48);
    \node at (-1.55,-1.43) {$(Z_2-1)_2$};
  \end{scope}
  % Vertex counterterm at the right end of the loop.
  \begin{scope}[xshift=-3.25cm,yshift=-1.45cm]
    \draw[line width=0.7pt] (-3.00,0) -- (-1.45,0);
    \draw[line width=0.7pt] (1.45,0) -- (3.00,0);
    \draw[electron] (-1.45,0) arc[start angle=180,end angle=0,x radius=1.45cm,y radius=0.88cm];
    \draw[photon] (1.45,0) arc[start angle=0,end angle=-180,x radius=1.45cm,y radius=0.88cm];
    \draw[line width=0.8pt] (1.33,-0.12)--(1.57,0.12);
    \draw[line width=0.8pt] (1.33,0.12)--(1.57,-0.12);
    \draw[callout] (1.95,-1.12) -- (1.62,-0.45);
    \node at (1.72,-1.42) {$(Z_2-1)_2$};
  \end{scope}
  % Electron-field counterterm.
  \begin{scope}[xshift=3.25cm,yshift=-1.45cm]
    \draw[line width=0.7pt] (-3.00,0) -- (3.00,0);
    \draw[line width=0.8pt] (-0.12,-0.13)--(0.12,0.13);
    \draw[line width=0.8pt] (-0.12,0.13)--(0.12,-0.13);
    \draw[callout] (-0.55,-1.10) -- (-0.12,-0.42);
    \node at (0,-1.42) {$(Z_3-1)_{\mathrm{overlap}}$};
  \end{scope}
\end{tikzpicture}
\caption{Some fourth-order graphs for the electron self-energy in quantum electrodynamics that involve overlapping divergences. Wavy lines are photons, other lines are electrons. The crosses mark the contribution of counterterms.}
\label{fig:weinberg-12-3}
\end{figure}
\endgroup
